\usepackage[a4paper,left=25mm,right=25mm,top=23mm,bottom=22mm,headheight=15pt]{geometry}
\usepackage{fontspec}
\usepackage{amsmath}
% TeX 自带的 Fandol 字体按文件名加载，由 kpathsea 定位字体目录。
\IfFontExistsTF{Microsoft YaHei}{\setCJKsansfont{Microsoft YaHei}}{\IfFontExistsTF{Noto Sans CJK SC}{\setCJKsansfont{Noto Sans CJK SC}}{\setCJKsansfont{FandolHei-Regular.otf}}}
\IfFontExistsTF{SimSun}{\setCJKmainfont{SimSun}[BoldFont=Microsoft YaHei]}{\IfFontExistsTF{Noto Serif CJK SC}{\setCJKmainfont{Noto Serif CJK SC}}{\setCJKmainfont{FandolSong-Regular.otf}[BoldFont=FandolHei-Regular.otf]}}
\IfFileExists{texgyrepagella-regular.otf}{\setmainfont{TeX Gyre Pagella}}{\setmainfont{Latin Modern Roman}}
\IfFileExists{texgyreheros-regular.otf}{\setsansfont{TeX Gyre Heros}}{\setsansfont{Latin Modern Sans}}
\usepackage{xcolor}
\definecolor{Ink}{HTML}{19333C}
\definecolor{Accent}{HTML}{BA643F}
\definecolor{Muted}{HTML}{63747A}
\definecolor{Paper}{HTML}{F5F2EA}
\usepackage{tikz}
\usepackage{fancyhdr}
\usepackage{hyperref}
\hypersetup{unicode=true,colorlinks=true,linkcolor=Ink,urlcolor=Accent,citecolor=Ink,pdftitle={分布式系统中的幂等性：从一次重复扣款开始},pdfsubject={幂等、事务与故障恢复}}
\hypersetup{pdfauthor={Allein},pdfsubject={分布式系统中的幂等性}}
\setlength{\parindent}{0pt}
\pagestyle{fancy}
\fancyhf{}
\fancyfoot[C]{\small\color{Muted}\thepage}
\renewcommand{\headrulewidth}{0pt}
\setcounter{tocdepth}{1}

\usepackage{graphicx}
\usepackage{caption}
\usetikzlibrary{arrows.meta,positioning,calc}
\captionsetup{font=small,labelfont={color=Accent},textfont={color=Ink},skip=7pt}
\setlength{\parindent}{2em}
\setlength{\parskip}{.35em}
\linespread{1.25}
\ctexset{chapter={format=\sffamily\bfseries\color{Ink}\LARGE,beforeskip=4pt,afterskip=20pt},section={format=\sffamily\bfseries\color{Ink}\large,beforeskip=14pt,afterskip=8pt}}
\tikzset{every picture/.style={font=\small\sffamily,text=Ink},flow/.style={-{Stealth[length=2mm]},draw=Ink,line width=.7pt},event/.style={fill=Ink,circle,inner sep=2pt}}

\clubpenalty=10000\widowpenalty=10000
\interfootnotelinepenalty=10000

% 正文中的代码用于解释机制，控制流程采用语言无关的伪代码。
\usepackage{fvextra}
\usepackage{tabularx}
\usepackage{booktabs}
\usepackage{placeins}
\IfFontExistsTF{Consolas}{\setmonofont{Consolas}}{\setmonofont{Latin Modern Mono}}
\IfFontExistsTF{SimSun}{\setCJKmonofont{SimSun}}{\IfFontExistsTF{Noto Serif CJK SC}{\setCJKmonofont{Noto Serif CJK SC}}{\setCJKmonofont{FandolSong-Regular.otf}}}
\DefineVerbatimEnvironment{BookCode}{Verbatim}{fontsize=\footnotesize,
  baselinestretch=1.04,breaklines=true,breakanywhere=true,samepage=true,
  xleftmargin=1em,xrightmargin=.5em,tabsize=4}
% 本章报文与伪代码均为短段，整段排入同一页，避免首行或请求头落单。
\BeforeBeginEnvironment{BookCode}{\par\addvspace{\medskipamount}\noindent\begin{minipage}{\linewidth}}
\AfterEndEnvironment{BookCode}{\end{minipage}\par\addvspace{\medskipamount}}
% 代码前可断行；末尾不强加断点，让 xeCJK 保持随后中文标点的禁则。
\newcommand{\code}[1]{\allowbreak{\ttfamily\def\_{\textunderscore\allowbreak}#1}}
\setcounter{secnumdepth}{1}
\ctexset{subsection={format=\sffamily\bfseries\color{Ink}\normalsize,beforeskip=13pt,afterskip=6pt}}
\tikzset{
  booknode/.style={draw=Ink!45,fill=Paper,rounded corners=2pt,
    align=flush center,inner sep=6pt,font=\small\sffamily},
  bookresult/.style={booknode,draw=Accent!65,fill=Accent!7},
  booknote/.style={font=\footnotesize\sffamily,text=Muted,align=flush center},
  bookarrow/.style={-{Stealth[length=2mm]},draw=Ink,line width=.7pt},
  bookalert/.style={-{Stealth[length=2mm]},draw=Accent,line width=.8pt}
}
\usetikzlibrary{fit,backgrounds,shapes.geometric}
\makeatletter
\setlength{\@fptop}{0pt}
\setlength{\@fpsep}{18pt}
\setlength{\@fpbot}{0pt plus 1fil}
\makeatother

% 中文引文使用明确的倾斜字形，西文沿用正文字族的斜体。
\IfFontExistsTF{SimSun}{%
  \newCJKfontfamily\BookQuoteCJK[ItalicFont=SimSun,ItalicFeatures={FakeSlant=.18}]{SimSun}%
}{\IfFontExistsTF{Noto Serif CJK SC}{%
  \newCJKfontfamily\BookQuoteCJK[ItalicFont={Noto Serif CJK SC},ItalicFeatures={FakeSlant=.18}]{Noto Serif CJK SC}%
}{%
  \newCJKfontfamily\BookQuoteCJK[ItalicFont=FandolSong-Regular.otf,ItalicFeatures={FakeSlant=.18}]{FandolSong-Regular.otf}%
}}
\newcommand{\DeclareBookEpigraph}[4]{%
  \expandafter\gdef\csname BookEpigraphText@#1\endcsname{#2}%
  \expandafter\gdef\csname BookEpigraphAuthor@#1\endcsname{#3}%
  \expandafter\gdef\csname BookEpigraphURL@#1\endcsname{#4}%
}
\newcommand{\BookEpigraph}[1]{%
  \par\vspace{5pt}
  \noindent\begin{minipage}{.90\linewidth}
    \setlength{\parindent}{0pt}\setlength{\parskip}{0pt}
    {\BookQuoteCJK\itshape\fontsize{10}{15}\selectfont\color{Ink}%
      \csname BookEpigraphText@#1\endcsname\par}
    \vspace{3pt}
    {\raggedleft\normalfont\fontsize{8.5}{12}\selectfont\color{Muted}%
      ——\href{\csname BookEpigraphURL@#1\endcsname}%
      {\csname BookEpigraphAuthor@#1\endcsname}\par}
  \end{minipage}\par\vspace{7pt}
}
% 章首引语及所据原文出处。
\DeclareBookEpigraph{preface}%
  {智慧是经验的女儿。}%
  {列奥纳多·达·芬奇，《笔记》§1150}%
  {https://www.gutenberg.org/ebooks/5000}

\DeclareBookEpigraph{introduction}%
  {世界是事实的总和，而非事物的总和。}%
  {路德维希·维特根斯坦，《逻辑哲学论》1.1}%
  {https://www.gutenberg.org/files/5740/5740-pdf.pdf}

\DeclareBookEpigraph{overview}%
  {观念的次序和联系，与事物的次序和联系相同。}%
  {巴鲁赫·斯宾诺莎，《伦理学》第二部分，命题七}%
  {https://www.gutenberg.org/ebooks/3800}

\DeclareBookEpigraph{01}%
  {程序必须写给人阅读，机器执行只是附带的目的。}%
  {阿贝尔森、杰拉尔德·萨斯曼与朱莉·萨斯曼，《计算机程序的构造和解释》初版序}%
  {https://sarabander.github.io/sicp/html/Preface-1e.xhtml}

\DeclareBookEpigraph{02}%
  {一切关于事实的推理，似乎都建立在因果关系之上。}%
  {大卫·休谟，《人类理解研究》第四节，§22}%
  {https://www.gutenberg.org/ebooks/9662}

\DeclareBookEpigraph{03}%
  {名称指谓对象，对象就是名称的意义。}%
  {路德维希·维特根斯坦，《逻辑哲学论》3.203}%
  {https://www.gutenberg.org/files/5740/5740-pdf.pdf}

\DeclareBookEpigraph{04}%
  {在分布式系统中，要认识到事件发生的次序只是一个偏序。}%
  {莱斯利·兰伯特，《时间、时钟与分布式系统中的事件排序》结论}%
  {https://lamport.azurewebsites.net/pubs/time-clocks.pdf}

\DeclareBookEpigraph{05}%
  {我们为自己构造事实的图像。}%
  {路德维希·维特根斯坦，《逻辑哲学论》2.1}%
  {https://www.gutenberg.org/files/5740/5740-pdf.pdf}

\DeclareBookEpigraph{06}%
  {程序测试可以用来显示错误的存在，却永远不能显示错误的不存在。}%
  {艾兹赫尔·戴克斯特拉，《结构化程序设计札记》EWD249}%
  {https://www.cs.utexas.edu/~EWD/transcriptions/EWD02xx/EWD249/EWD249.html}

\DeclareBookEpigraph{part1}%
  {从最简单、最容易认识的对象开始，一点一点、逐步上升，直到认识较为复杂的对象。}%
  {勒内·笛卡尔，《方法谈》第二部分，第三条准则}%
  {https://www.gutenberg.org/ebooks/59}

\DeclareBookEpigraph{07}%
  {把任务分开是好事，同时，我们还得把分开后的端头重新接起来。}%
  {艾兹赫尔·戴克斯特拉，《结构化程序设计札记》EWD249}%
  {https://www.cs.utexas.edu/~EWD/transcriptions/EWD02xx/EWD249/EWD249.html}

\DeclareBookEpigraph{08}%
  {每个物体都保持静止或匀速直线运动的状态，除非外加的力迫使它改变这种状态。}%
  {艾萨克·牛顿，《自然哲学的数学原理》运动第一定律}%
  {https://www.gutenberg.org/ebooks/28233}

\DeclareBookEpigraph{09}%
  {通信的根本问题，是在一个地点精确或近似地重现另一个地点选定的消息。}%
  {克劳德·香农，《通信的数学理论》引言}%
  {https://people.math.harvard.edu/~ctm/home/text/others/shannon/entropy/entropy.pdf}

\DeclareBookEpigraph{10}%
  {我们通过确立约定的接口来控制复杂性，使标准而充分理解的部件能够自由组合，构成系统。}%
  {阿贝尔森、杰拉尔德·萨斯曼与朱莉·萨斯曼，《计算机程序的构造和解释》初版序}%
  {https://sarabander.github.io/sicp/html/Preface-1e.xhtml}

\DeclareBookEpigraph{part2}%
  {整体大于部分。}%
  {欧几里得，《几何原本》第一卷，公理五}%
  {https://mathcs.clarku.edu/~djoyce/elements/bookI/cn.html}

\DeclareBookEpigraph{11}%
  {现在孕育着未来。}%
  {戈特弗里德·威廉·莱布尼茨，《单子论》§22}%
  {https://www.gutenberg.org/ebooks/17641}

\DeclareBookEpigraph{12}%
  {好的记号使头脑摆脱一切不必要的工作，能够集中精力处理更高深的问题。}%
  {阿尔弗雷德·诺思·怀特海，《数学导论》第五章}%
  {https://www.gutenberg.org/files/41568/41568-pdf.pdf}

\DeclareBookEpigraph{13}%
  {数学是一门给不同事物以相同名称的艺术。这些事物虽在内容上不同，在形式上却应当相似。}%
  {亨利·庞加莱，《科学与方法》第一编第二章}%
  {https://www.gutenberg.org/ebooks/39713}

\DeclareBookEpigraph{14}%
  {我们应当把宇宙的现在状态，看作先前状态的结果，以及随后状态的原因。}%
  {皮埃尔-西蒙·拉普拉斯，《概率的哲学论》第二章}%
  {https://www.gutenberg.org/ebooks/58881}

\DeclareBookEpigraph{15}%
  {因此，明智的人会使自己的信念与证据相称。}%
  {大卫·休谟，《人类理解研究》第十节，§87}%
  {https://www.gutenberg.org/ebooks/9662}

\DeclareBookEpigraph{part3}%
  {真理是时间唯一的女儿。}%
  {列奥纳多·达·芬奇，《笔记》§1152}%
  {https://www.gutenberg.org/ebooks/5000}

\DeclareBookEpigraph{16}%
  {把所考察的每一个困难，尽可能分成充分解决它所需要的若干部分。}%
  {勒内·笛卡尔，《方法谈》第二部分，第二条准则}%
  {https://www.gutenberg.org/ebooks/59}

\DeclareBookEpigraph{17}%
  {我们通过构造抽象，在适当的时候隐藏细节，以此控制复杂性。}%
  {阿贝尔森、杰拉尔德·萨斯曼与朱莉·萨斯曼，《计算机程序的构造和解释》初版序}%
  {https://sarabander.github.io/sicp/html/Preface-1e.xhtml}

\DeclareBookEpigraph{18}%
  {形式是这种内在意义的外在表达。}%
  {瓦西里·康定斯基，《论艺术的精神》第六章}%
  {https://www.gutenberg.org/ebooks/5321}

\DeclareBookEpigraph{19}%
  {如同在数学中一样，在自然哲学中，用分析方法探究困难问题，应当先于综合方法。}%
  {艾萨克·牛顿，《光学》第三卷，疑问三十一}%
  {https://www.gutenberg.org/ebooks/33504}

\DeclareBookEpigraph{20}%
  {文明的进步，在于扩大那些我们无须思索便能完成的重要操作的数量。}%
  {阿尔弗雷德·诺思·怀特海，《数学导论》第五章}%
  {https://www.gutenberg.org/files/41568/41568-pdf.pdf}

\DeclareBookEpigraph{21}%
  {我们只能看见前方不远的地方，但已能看见那里有许多工作要做。}%
  {艾伦·图灵，《计算机器与智能》结尾}%
  {https://www.csee.umbc.edu/courses/471/papers/turing.pdf}

\DeclareBookEpigraph{part4}%
  {每一件艺术作品都是它所属时代的孩子，而且在许多情况下，也是我们情感的母亲。}%
  {瓦西里·康定斯基，《论艺术的精神》第一章}%
  {https://www.gutenberg.org/ebooks/5321}

\DeclareBookEpigraph{closing}%
  {科学由事实构成，如同房屋由石头建成；一堆事实还构不成科学，正如一堆石头还构不成房屋。}%
  {亨利·庞加莱，《科学与假设》第九章}%
  {https://www.gutenberg.org/ebooks/39713}

